\clearpage
\section{Question 2}

\begin{colbox}
	\begin{enumerate}
		\item Write down the Hamiltonian for the Neutron. Write down the time-dependent Schrödinger equation for the state of the Neutron and show it can be decomposed into two similar equations for $\psi_\pm$
		\item Show that if $\ket{\psi(t)}$ is of the form in the previous question then 
		\begin{equation*}
			\diff*{\ins{\int\left\lvert\psi_\pm(\vec r, t)\right\rvert\dd \vec r}}{t} = 0
		\end{equation*} What can we deduce for the probability to measure $\mu_z = \pm\mu_0$
		\item Long question I won't copy
		\item Use Ehrenfest's theorem to find $\diff*{\left\langle\vec r\right\rangle_\pm}{t}$ and $\diff*{\left\langle\vec p\right\rangle_\pm}{t}$. Solve the resulting system of equations and find the time development for both. Give a physical interpretation for the result and explain why the original beam splits into two beams.
		\item What is the distance D between the two marks on the screen when $L=1m$? Assume $L\gg D$. What is the ratio of the lengths $D$ and the Neutron's radius?
	\end{enumerate}
\end{colbox}

\subsection{Hamiltonian}

As we know the Hamiltonian is

\begin{equation}
	H = \frac{p^2}{2m}-\vec\mu\cdot\vec B = \frac{p^2}{2m}\mathds{1}-\mu_0\ins{B_0+bz}\sigma_z
\end{equation}

And the time-dependent Schrödinger equation is

\begin{equation}
	i\hbar \diff*{\ket{\psi(t)}}{t} = H\ket{\psi(t)}
\end{equation}

If we replace the generic wave equation with the vector wave equation from question 1 then the ODE that is the Schrödinger equation becomes two:

\begin{equation}
	i\hbar \diff*{\psi(\vec r,t)\pmat{\alpha_+\\\alpha_-}}{t} = H_\pm\ket{\psi(t)} = -\frac{\hbar^2}{2m}\nabla^2\psi_\pm\mp\mu_0\ins{B_0+bz}\psi_\pm
\end{equation}

The $\mp$ appearing because of the inverted magnetic moment.

\subsection{Integral Derivative}

Since $\ket{\psi(t)}$ is of the form $\ket{\psi(\vec r,t)}=\psi(\vec r,t)\pmat{\alpha_+\\\alpha_-}$, then 

\begin{equation}
	\int\left\lvert\psi_\pm(\vec r,t)\right\rvert\dd \vec r = \left\lvert\alpha_\pm\right\rvert^2
\end{equation}

then it is a constant and therefore also not time-dependent, which gives us the needed identity. This expression is also the same as the probability of measuring $\mu_\pm$, which means the probabilities of those are also constant.

\subsection{Part 3}

From the structure of the integral we can tell that this is some kind of observation value. The values noted represent the probability of finding the particle in a certain position.

\subsection{Time Development}

Recalling Ehrenfest's Theorem:

\begin{equation}
	i\hbar \diff*{\langle\vec\mu\rangle}{t} = \left\langle[H,\vec\mu]\right\rangle
\end{equation}

then (shortening $\vec p_\pm$ to $p$)

\begin{align}
	\diff*{\left\langle\vec r_\pm\right\rangle}{t} &= -\frac{i}{\hbar}\left\langle[H,\vec r]\right\rangle\\
	&= -\frac{i}{\hbar} \int\psi^*[H,\vec r]\psi\dd \vec r\\
	&= -\frac{i}{\hbar} \int\psi^*[\frac{p^2}{2m},\vec r]\psi\dd \vec r\\
	&= -\frac{i}{2m\hbar} \int\psi^*[p^2,\vec r]\psi\dd \vec r
\end{align}

where

\begin{align}
	[p^2,r] &= p^2r-rp^2\\
	&= ppr-rpp\\
	&= p(rp-(rp-pr))-(pr-(pr-rp))p\\
	&= p(rp-[r,p])-(pr-[p,r])p\\
	&= prp-p[r,p]-prp+[p,r]p\\
	&= 2i\hbar p
\end{align}

therefore

\begin{align}
	-\frac{i}{2m\hbar} \int\psi^*[p^2,\vec r]\psi\dd \vec r &= \frac{1}{m}\int\psi^*p\psi\dd \vec r\\
	&= \frac{\langle \vec p_\pm\rangle}{m}
\end{align}

and

\begin{align}
	\diff*{\langle\vec p_\pm\rangle}{t} &= -\frac{i}{\hbar}\langle[H,\vec p_\pm]\rangle\\
\end{align}

we notice immediately that 

\begin{align}
	\diff*{\langle p_x\rangle}{t} = \diff*{\langle p_y\rangle}{t} = 0
\end{align}

and now

\begin{align}
	\diff*{\langle p_{z\pm}\rangle}{t} &= -\frac{i}{\hbar}\int\psi^*[H,p_{z\pm}]\psi\dd\vec r\\
	&= -\frac{i}{\hbar}\int\psi^*[-\mu_0bz,p_{z\pm}]\psi\dd\vec r\\
	&= \frac{i\mu_0b}{\hbar}\int\psi^*[z,p_{z\pm}]\psi\dd \vec r\\
	&= -\mp\mu_0b\int \psi^*\psi\dd \vec r\\
	&= \pm\mu_0b
\end{align}

Therefore

\begin{align}
	\vec p_\pm = \pmat{p_0\\0\\\pm\mu_0bt} && \vec r = \pmat{0\\0\\\frac{\pm\mu_0bt^2}{2m}}
\end{align}

\subsection{Finding D}

\begin{align}
	D = \langle z_+ \rangle - \langle z_- \rangle = \frac{\mu_0bt^2}{m} = \frac{\mu_0bL^2}{mv^2} &= \frac{\mu_0bL^2}{2E}\\
	= \frac{\mu_0bL^2}{2\ins{\frac{p^2}{2m}}} = \frac{m\mu_0bL^2}{p^2}&=\frac{m\mu_0bL^2\lambda^2}{h^2}\\
	= [\text{inserting the numbers into a calculator}] &= \SI{2.06e-5}{\meter}
\end{align}
which gives us
\begin{equation}
	\frac{D}{r_n} = \frac{\SI{2.06e-5}{\meter}}{\SI{0.8e-15}{\meter}}=\SI{2.575e10}{}
\end{equation}